% =====================================================================
% RASCUNHO (2026-11-06): discussão de vazão e atraso, para a autora trabalhar.
%   Lugar: depois dos resultados de hardware (intervalo, área e potência), que ainda não
%   estão escritos. Quando eles entrarem, mover esta seção para depois deles.
% DADOS: hardware/resultados_servidor/atraso_por_bloco.xlsx (fórmulas) e
%   latencia_por_bloco.csv, gerados por hardware/latencia_por_bloco.py e
%   hardware/atraso_por_bloco_xlsx.py a partir dos logs xsim do servidor (Vivado 2023.2).
%   Os logs têm só 3 lotes (ou 20 matrizes); o intervalo sai de 2 (ou 19) medidas.
% A tabela mostra só <24,12>. Os outros formatos estão na planilha:
%   - a 455 passa de T_c de +3,2 % (<18,8>) a +11,5 % (<30,14>);
%   - a 152 passa de T_c por cerca de 21,5 % em todos;
%   - a folga da 606 cai de 11 % (<18,8>) para 3 % (<30,14>).
% =====================================================================
\section{\texorpdfstring{\textcolor{red}{Vazão e atraso}}{Vazão e atraso}}\label{sec:resultados-atraso}

\textcolor{red}{As arquiteturas foram projetadas para atender o limite de vazão de cada condição de teste (Seção~\ref{subsec:restricao-vazao}). Esta seção confere esse limite e mostra o atraso que resulta dessa escolha.}

% ---------------------------------------------------------------------
\subsection{\texorpdfstring{\textcolor{red}{Vazão}}{Vazão}}
% ---------------------------------------------------------------------
\textcolor{red}{A cada $T_c$, chega um bloco de $N_{\text{PRB}} = 275$ matrizes $\mathbf G_c$, uma por PRB. As arquiteturas 606, 455, 152 e 114 processam as matrizes em lotes de 12, 16, 32 e 32 matrizes. Um lote novo entra no módulo a cada intervalo de lote, antes de o anterior terminar. O tempo para o módulo aceitar todo o bloco é o número de lotes vezes o intervalo de lote. Em todas as arquiteturas e em todos os formatos, esse tempo fica abaixo de $T_c$. Assim, o módulo acompanha o ritmo de chegada das matrizes, e não se forma fila.}
% FEITO (2026-11-06): "lote" definido em subsubsec:hw-lotes; o fim do bloco é a Eq. eq:fim-tarefas.
% ANTIGO: % TODO: definir "lote" na Metodologia, junto da Tabela tab:hw-arquiteturas.

% ---------------------------------------------------------------------
\subsection{\texorpdfstring{\textcolor{red}{Atraso}}{Atraso}}
% ---------------------------------------------------------------------
\textcolor{red}{O atraso vai da chegada da primeira amostra de um bloco até a última inversão desse bloco. Ele tem duas partes. A primeira é o tempo de coleta, que o receptor espera para juntar as $N$ amostras (Seção~\ref{subsec:pipeline}). Com $N$ amostras e 48 REs de DM-RS por PRB e por slot, a coleta leva $\lceil N/48 \rceil$ slots, limitada a $T_c$. A segunda é o tempo para inverter o bloco: a latência do primeiro lote mais o intervalo de cada lote seguinte. A Tabela~\ref{tab:atraso-bloco} mostra essas duas partes no formato $\langle 24, 12\rangle$.}

\begin{table}[!ht]
\caption{\textcolor{red}{Throughput and delay of one block of $N_{\text{PRB}} = 275$ inversions, format $\langle 24, 12\rangle$.}}
\label{tab:atraso-bloco}
\vspace{4pt} % a classe zera o espaço entre a caption e a tabela
\centering
\footnotesize
\color{red}
\renewcommand{\arraystretch}{1.2}
\setlength{\tabcolsep}{4pt}
\begin{tabular}{llccccc}
\toprule
\textbf{Architecture} & \begin{tabular}[c]{@{}c@{}}\textbf{Test}\\\textbf{condition}\end{tabular} & \begin{tabular}[c]{@{}c@{}}{\boldmath $T_c$}\\\textbf{(ms)}\end{tabular} & \begin{tabular}[c]{@{}c@{}}\textbf{Block input}\\\textbf{(ms)}\end{tabular} & \begin{tabular}[c]{@{}c@{}}\textbf{Block end}\\\textbf{(ms)}\end{tabular} & \begin{tabular}[c]{@{}c@{}}\textbf{Collection}\\\textbf{(ms)}\end{tabular} & \begin{tabular}[c]{@{}c@{}}\textbf{Total delay}\\\textbf{(ms)}\end{tabular} \\
\midrule
9090 & TDLA30-5    & 25.000 & 20.884 & 20.884 & 2.000 & 22.884 \\
4545 & TDLA30-10   & 12.500 & 8.305  & 8.354  & 2.000 & 10.354 \\
606  & TDLA30-75   & 1.667  & 1.391  & 1.543  & 0.500 & 2.043 \\
455  & TDLC300-100 & 1.250  & 1.159  & 1.336  & 1.250 & 2.586 \\
152  & TDLA30-300  & 0.417  & 0.369  & 0.506  & 0.417 & 0.923 \\
114  & TDLB100-400 & 0.313  & 0.184  & 0.265  & 0.313 & 0.578 \\
\bottomrule
\end{tabular}
\vspace{4pt}
\legend{\textcolor{red}{Source: the author, from the netlist simulation at 100~MHz. Block input: time to accept the 275 matrices. Block end: time until the last inversion of the block. Total delay: collection plus block end.}}
\end{table}

\textcolor{red}{Nas arquiteturas 455 e 152, a última inversão do bloco termina depois de $T_c$. Esse excesso não se acumula, pois o bloco seguinte já entra no módulo a tempo. Ele é um atraso fixo, de cerca de uma latência de lote, e não uma perda de vazão. Além disso, cada PRB só precisa da sua própria inversão. Os primeiros PRBs ficam prontos cerca de uma latência de lote depois da chegada do bloco. O atraso da tabela é o do último PRB.}

\textcolor{red}{A ITU-R prevê que o IMT-2020 possa oferecer latência de 1~ms pela interface aérea \cite[p.~14]{itu-m2083}. Essa latência é a contribuição da rede de rádio para o tempo entre o envio e a recepção de um pacote \cite[p.~13]{itu-m2083}. O processamento na BS é uma das parcelas da latência do plano de usuário \cite[Sec.~7.2.6 e Tab.~3]{itu-m2412}. Só as arquiteturas 152 e 114 ficam abaixo de 1~ms. Nas outras, o atraso vai de 2,0 a 22,9~ms.}

\textcolor{red}{Parte desse atraso vem do método de estimação, e não do hardware. Com $N = 192$ em FR1, a coleta leva 4 slots de 0,5~ms, ou seja, 2~ms. Esse tempo já passa de 1~ms antes de qualquer inversão. A outra parte vem da escolha de projeto. As arquiteturas das condições de baixa mobilidade distribuem as 275 inversões por quase todo o $T_c$. Com isso, elas gastam menos área e potência, mas o último PRB fica pronto mais tarde. Uma arquitetura mais rápida reduziria essa parte do atraso, com mais área e potência. A arquitetura 114, por exemplo, termina um bloco em 0,27~ms.}
% TODO: comparar com a área e a potência das arquiteturas (resultados de hardware ainda não escritos).

\textcolor{red}{A coleta também depende do PUSCH. O DM-RS só é transmitido junto com o PUSCH, e o PUSCH só ocupa os slots em que os UEs alvo foram escalonados (Seção~\ref{subsec:simbolos}). Este trabalho supõe que todo slot leva PUSCH. As bandas n78 e n257 usam TDD \cite[Tabs.~5.2-1 e 5.2-2]{3gpp-ts38104}, e só os slots de \textit{uplink} levam o PUSCH. Com um padrão TDD real, os slots de \textit{uplink} ficam mais espaçados. A coleta das $N$ amostras leva mais tempo ou, se passar de $T_c$, o número de amostras disponível diminui.}
% TODO: padrão TDD típico (sem fonte local). Se for usado um número, buscar a fonte.
% TODO: o valor de 4 ms para eMBB está no Rep. ITU-R M.2410-0, que não está em books/.
